\documentclass[a4paper,12pt,twoside]{memoir}

% Castellano
\usepackage[spanish,es-tabla]{babel}
\selectlanguage{spanish}
\usepackage[utf8]{inputenc}
\usepackage[T1]{fontenc}
\usepackage{lmodern} % Scalable font
\usepackage{microtype}
\usepackage{placeins}

\RequirePackage{booktabs}
\RequirePackage[table]{xcolor}
\RequirePackage{xtab}
\RequirePackage{multirow}

% Links
\PassOptionsToPackage{hyphens}{url}\usepackage[colorlinks]{hyperref}
\hypersetup{
	allcolors = {blue}
}

\usepackage[
  backend=biber,
  style=numeric
]{biblatex}
\addbibresource{bibliografia.bib}

% Ecuaciones
\usepackage{amsmath}

% Rutas de fichero / paquete
\newcommand{\ruta}[1]{{\sffamily #1}}

% Párrafos
\nonzeroparskip

% Huérfanas y viudas
\widowpenalty100000
\clubpenalty100000


% Comando para insertar una imagen en un lugar concreto.
% Los parámetros son:
% 1 --> Ruta absoluta/relativa de la figura
% 2 --> Texto a pie de figura
% 3 --> Tamaño en tanto por uno relativo al ancho de página
\usepackage{graphicx}
\newcommand{\imagen}[3]{
	\begin{figure}[!h]
		\centering
		\includegraphics[width=#3\textwidth]{#1}
		\caption{#2}\label{fig:#1}
	\end{figure}
	\FloatBarrier
}

% Comando para insertar una imagen sin posición.
% Los parámetros son:
% 1 --> Ruta absoluta/relativa de la figura
% 2 --> Texto a pie de figura
% 3 --> Tamaño en tanto por uno relativo al ancho de página
\newcommand{\imagenflotante}[3]{
	\begin{figure}
		\centering
		\includegraphics[width=#3\textwidth]{#1}
		\caption{#2}\label{fig:#1}
	\end{figure}
}


% El comando \figura nos permite insertar figuras comodamente, y utilizando
% siempre el mismo formato. Los parametros son:
% 1 --> Porcentaje del ancho de página que ocupará la figura (de 0 a 1)
% 2 --> Fichero de la imagen
% 3 --> Texto a pie de imagen
% 4 --> Etiqueta (label) para referencias
% 5 --> Opciones que queramos pasarle al \includegraphics
% 6 --> Opciones de posicionamiento a pasarle a \begin{figure}
\newcommand{\figuraConPosicion}[6]{%
  \setlength{\anchoFloat}{#1\textwidth}%
  \addtolength{\anchoFloat}{-4\fboxsep}%
  \setlength{\anchoFigura}{\anchoFloat}%
  \begin{figure}[#6]
    \begin{center}%
      \Ovalbox{%
        \begin{minipage}{\anchoFloat}%
          \begin{center}%
            \includegraphics[width=\anchoFigura,#5]{#2}%
            \caption{#3}%
            \label{#4}%
          \end{center}%
        \end{minipage}
      }%
    \end{center}%
  \end{figure}%
}

%
% Comando para incluir imágenes en formato apaisado (sin marco).
\newcommand{\figuraApaisadaSinMarco}[5]{%
  \begin{figure}%
    \begin{center}%
    \includegraphics[angle=90,height=#1\textheight,#5]{#2}%
    \caption{#3}%
    \label{#4}%
    \end{center}%
  \end{figure}%
}
% Para las tablas
\newcommand{\otoprule}{\midrule [\heavyrulewidth]}
%
% Nuevo comando para tablas pequeñas (menos de una página).
\newcommand{\tablaSmall}[5]{%
 \begin{table}
  \begin{center}
   \rowcolors {2}{gray!35}{}
   \begin{tabular}{#2}
    \toprule
    #4
    \otoprule
    #5
    \bottomrule
   \end{tabular}
   \caption{#1}
   \label{tabla:#3}
  \end{center}
 \end{table}
}

%
% Nuevo comando para tablas pequeñas (menos de una página).
\newcommand{\tablaSmallSinColores}[5]{%
 \begin{table}[H]
  \begin{center}
   \begin{tabular}{#2}
    \toprule
    #4
    \otoprule
    #5
    \bottomrule
   \end{tabular}
   \caption{#1}
   \label{tabla:#3}
  \end{center}
 \end{table}
}

\newcommand{\tablaApaisadaSmall}[5]{%
\begin{landscape}
  \begin{table}
   \begin{center}
    \rowcolors {2}{gray!35}{}
    \begin{tabular}{#2}
     \toprule
     #4
     \otoprule
     #5
     \bottomrule
    \end{tabular}
    \caption{#1}
    \label{tabla:#3}
   \end{center}
  \end{table}
\end{landscape}
}

%
% Nuevo comando para tablas grandes con cabecera y filas alternas coloreadas en gris.
\newcommand{\tabla}[6]{%
  \begin{center}
    \tablefirsthead{
      \toprule
      #5
      \otoprule
    }
    \tablehead{
      \multicolumn{#3}{l}{\small\sl continúa desde la página anterior}\\
      \toprule
      #5
      \otoprule
    }
    \tabletail{
      \hline
      \multicolumn{#3}{r}{\small\sl continúa en la página siguiente}\\
    }
    \tablelasttail{
      \hline
    }
    \bottomcaption{#1}
    \rowcolors {2}{gray!35}{}
    \begin{xtabular}{#2}
      #6
      \bottomrule
    \end{xtabular}
    \label{tabla:#4}
  \end{center}
}

%
% Nuevo comando para tablas grandes con cabecera.
\newcommand{\tablaSinColores}[6]{%
  \begin{center}
    \tablefirsthead{
      \toprule
      #5
      \otoprule
    }
    \tablehead{
      \multicolumn{#3}{l}{\small\sl continúa desde la página anterior}\\
      \toprule
      #5
      \otoprule
    }
    \tabletail{
      \hline
      \multicolumn{#3}{r}{\small\sl continúa en la página siguiente}\\
    }
    \tablelasttail{
      \hline
    }
    \bottomcaption{#1}
    \begin{xtabular}{#2}
      #6
      \bottomrule
    \end{xtabular}
    \label{tabla:#4}
  \end{center}
}

%
% Nuevo comando para tablas grandes sin cabecera.
\newcommand{\tablaSinCabecera}[5]{%
  \begin{center}
    \tablefirsthead{
      \toprule
    }
    \tablehead{
      \multicolumn{#3}{l}{\small\sl continúa desde la página anterior}\\
      \hline
    }
    \tabletail{
      \hline
      \multicolumn{#3}{r}{\small\sl continúa en la página siguiente}\\
    }
    \tablelasttail{
      \hline
    }
    \bottomcaption{#1}
  \begin{xtabular}{#2}
    #5
   \bottomrule
  \end{xtabular}
  \label{tabla:#4}
  \end{center}
}



\definecolor{cgoLight}{HTML}{EEEEEE}
\definecolor{cgoExtralight}{HTML}{FFFFFF}

%
% Nuevo comando para tablas grandes sin cabecera.
\newcommand{\tablaSinCabeceraConBandas}[5]{%
  \begin{center}
    \tablefirsthead{
      \toprule
    }
    \tablehead{
      \multicolumn{#3}{l}{\small\sl continúa desde la página anterior}\\
      \hline
    }
    \tabletail{
      \hline
      \multicolumn{#3}{r}{\small\sl continúa en la página siguiente}\\
    }
    \tablelasttail{
      \hline
    }
    \bottomcaption{#1}
    \rowcolors[]{1}{cgoExtralight}{cgoLight}

  \begin{xtabular}{#2}
    #5
   \bottomrule
  \end{xtabular}
  \label{tabla:#4}
  \end{center}
}



\graphicspath{ {img/} }

% Capítulos
\chapterstyle{bianchi}
\newcommand{\capitulo}[2]{
	\setcounter{chapter}{#1}
	\setcounter{section}{0}
	\setcounter{figure}{0}
	\setcounter{table}{0}
	\chapter*{#2}
	\addcontentsline{toc}{chapter}{#2}
	\markboth{#2}{#2}
}

% Apéndices
\renewcommand{\appendixname}{Apéndice}
\renewcommand*\cftappendixname{\appendixname}

\newcommand{\apendice}[1]{
	%\renewcommand{\thechapter}{A}
	\chapter{#1}
}

\renewcommand*\cftappendixname{\appendixname\ }

% Formato de portada
\makeatletter
\usepackage{xcolor}
\newcommand{\tutor}[1]{\def\@tutor{#1}}
\newcommand{\course}[1]{\def\@course{#1}}
\definecolor{cpardoBox}{HTML}{E6E6FF}
\def\maketitle{
  \null
  \thispagestyle{empty}
  % Cabecera ----------------
\noindent\includegraphics[width=\textwidth]{cabecera}\vspace{1cm}%
  \vfill
  % Título proyecto y escudo informática ----------------
  \colorbox{cpardoBox}{%
    \begin{minipage}{.8\textwidth}
      \vspace{.5cm}\Large
      \begin{center}
      \textbf{TFG del Grado en Ingeniería Informática}\vspace{.6cm}\\
      \textbf{\LARGE\@title{}}
      \end{center}
      \vspace{.2cm}
    \end{minipage}

  }%
  \hfill\begin{minipage}{.20\textwidth}
    \includegraphics[width=\textwidth]{escudoInfor}
  \end{minipage}
  \vfill
  % Datos de alumno, curso y tutores ------------------
  \begin{center}%
  {%
    \noindent\LARGE
    Presentado por \@author{}\\ 
    en Universidad de Burgos --- \@date{}\\
    Tutor: \@tutor{}\\
  }%
  \end{center}%
  \null
  \cleardoublepage
  }
\makeatother

\newcommand{\nombre}{Ignacio Zaldo González} %%% cambio de comando

% Datos de portada
\title{Liga de Tenis de los Colegios de Burgos}
\author{\nombre}
\tutor{César Represa Pérez}
\date{\today}

\begin{document}

\maketitle


\newpage\null\thispagestyle{empty}\newpage


%%%%%%%%%%%%%%%%%%%%%%%%%%%%%%%%%%%%%%%%%%%%%%%%%%%%%%%%%%%%%%%%%%%%%%%%%%%%%%%%%%%%%%%%
\thispagestyle{empty}


\noindent\includegraphics[width=\textwidth]{cabecera}\vspace{1cm}

\noindent D. César Represa Pérez, profesor del departamento de Ingeniería Electromecánica, área de Tecnología Electrónica.

\noindent Expone:

\noindent Que el alumno D. \nombre, con DNI 71707258C, ha realizado el Trabajo final de Grado en Ingeniería Informática titulado Liga de Tenis de los Colegios de Burgos. 

\noindent Y que dicho trabajo ha sido realizado por el alumno bajo la dirección del que suscribe, en virtud de lo cual se autoriza su presentación y defensa.

\begin{center} %\large
En Burgos, {\large \today}
\end{center}

\vfill\vfill\vfill

% Author and supervisor
\begin{minipage}{0.45\textwidth}
\begin{flushleft} %\large
Vº. Bº. del Tutor:\\[2cm]
D. nombre tutor
\end{flushleft}
\end{minipage}
\hfill
\begin{minipage}{0.45\textwidth}
\begin{flushleft} %\large
Vº. Bº. del co-tutor:\\[2cm]
D. nombre co-tutor
\end{flushleft}
\end{minipage}
\hfill

\vfill

% para casos con solo un tutor comentar lo anterior
% y descomentar lo siguiente
%Vº. Bº. del Tutor:\\[2cm]
%D. nombre tutor


\newpage\null\thispagestyle{empty}\newpage




\frontmatter

% Abstract en castellano
\renewcommand*\abstractname{Resumen}
\begin{abstract}
Este proyecto consiste en el desarrollo de una aplicación web dedicada a una liga de tenis para los colegios de Burgos, la cual permite tener un CRUD (Create, Read, Update, Delete -- Crear, Leer, Actualizar y Borrar) y ver información relevante sobre la liga (Calendario, clasificación, jugadores, colegios, partidos) actualizando datos y teniéndolos guardados en base de datos a través de APIs creadas por mí. Además, se muestran mapas con ubicaciones marcadas donde se van a jugar los partidos utilizando Leaflet y OpenStreetMap. 

También se predicen los resultados de los partidos que se jugarán más adelante basándose en resultados anteriores y estadísticas que se sacan del jugador que representará al colegio, esto se hace con la integración de otra API externa, la de Google AI Studio, que además de lo anterior, permite chatear con el bot para preguntarle información relevante sobre la liga.

Esta aplicación web está desarrollada en fullstack, es decir, hay un front-end hecho en Angular, que incluye HTML, CSS y TypeScript y un back-end que está hecho con .NET para crear las APIs y poder actualizar la información en las tablas de la base de datos hecha con HeidiSQL utilizando los servicios de MySQL.

\end{abstract}

\renewcommand*\abstractname{Descriptores}
\begin{abstract}
Liga de tenis, aplicación Web, Angular, .NET, bases de datos, estadísticas, CRUD, APIs, mapas, predicción, google ai, \ldots
\end{abstract}

\clearpage

% Abstract en inglés
\renewcommand*\abstractname{Abstract}
\begin{abstract}
This project involves the development of a full-stack web application for a school tennis league in Burgos. It features a CRUD interface (Create, Read, Update and Delete) and displays all relevant league data—schedule, standings, players, schools and matches—by updating and storing information in a MySQL database via custom APIs I’ve built.

Maps are integrated using Leaflet and OpenStreetMap to show the exact locations where matches will be played. The application also predicts upcoming match results based on past outcomes and player statistics, this can be possible because the app uses another external API, the Google AI Studio one, that allows doing that and also chatting with the bot in order to obtain relevant information about the league.

On the front end, the app is built with Angular (HTML, CSS and TypeScript), while the back end is implemented in .NET, exposing the APIs that interact with the HeidiSQL-managed MySQL database.
\end{abstract}

\renewcommand*\abstractname{Keywords}
\begin{abstract}
Tennis league, web application, Angular, .NET, databases, statistics, CRUD, APIs, maps, predictions, ai usage integrated with the web, \ldots
\end{abstract}

\clearpage

% Indices
\tableofcontents

\clearpage

\listoffigures

\clearpage

\listoftables
\clearpage

\mainmatter
\capitulo{1}{Introducción}

Durante varios años, el seguimiento del rendimiento deportivo a nivel de deportes escolares se ha basado en la utilización de métodos manuales como por ejemplo tablas con la clasificación en tablones de anuncios o en forma de papel. Sin embargo, con el auge de las tecnologías y la fácil accesibilidad a dispositivos que permiten ver páginas web por parte de todo el mundo, surge una gran oportunidad para digitalizar y optimizar la gestión y organización de estas ligas y hacer que todo esto sea más fácil de ver y gestionar.

En el mundo del tenis normalmente las competiciones que hay son torneos de corta duración en los que cada jugador va pasando por diferentes fases y eliminatorias hasta llegar a la final y tener opción a ganarla, pero en ese tipo de torneos se suele competir representando a la misma persona, es decir, cada jugador compite por él mismo y no por un equipo.

En este caso el planteamiento es algo diferente ya que la forma de competir será diferente a la clásica que se lleva a cabo en el tenis más profesional. Como es una liga entre colegios, su duración vendrá dada por el calendario escolar y no será en forma de torneo con una duración más corta. 

En esta liga se podrán inscribir los colegios que lo soliciten, pudiendo después cada uno de ellos inscribir a los jugadores que quieran. Los colegios competirán durante toda la temporada por ver cual es el que más puntos logra conseguir, asignando para cada partido el jugador que cada colegio elija.

Esta forma de plantear la competición permite que al jugar muchos partidos, jueguen muchos jugadores diferentes entre ellos y que aunque se enfrenten en un partido los mismos colegios que se enfrentaron en otro partido anterior, el resultado pueda ser muy diferente al poder ser otros jugadores los que se enfrenten en representación de ese colegio.

En una liga de tenis como esta, tenemos una competición en la que se inscriben colegios y juegan partidos entre ellos, aportando la victoria de un partido 2 puntos para la clasificación y la derrota 1 punto, al ser tenis, no existe el empate y por tanto se puntuará de esta manera.

Este proyecto se centra en el desarrollo de una aplicación web para la gestión de una liga de tenis escolar en Burgos. A diferencia de las plataformas tradicionales que solo permiten mostrar resultados, con esta aplicación se busca ofrecer una experiencia completa que combina la información, análisis y funcionalidades nuevas tanto para jugadores, como entrenadores, organizadores y público en general.

\imagen{mi-web.png}{Imagen de la página de inicio de la web}{1.0}

Este proyecto propone una solución innovadora utilizando tecnologías más modernas como angular para el front-end y .net para el back-end junto con una base de datos que permite almacenar y consultar información clave sobre los partidos, jugadores, colegios y la clasificación en una tabla global por cada colegio. La aplicación también contempla la integración de herramientas para facilitar el registro de resultados visualización de calendarios y estadísticas y un sistema de predicción de resultados basados en los datos históricos.

Uno de los elementos más destacados de la web es su enfoque en la accesibilidad y el dinamismo, permitiendo que cualquier usuario pueda seguir el desarrollo de la Liga en tiempo real y conocer qué jugadores son los mejores y descubrir cómo pueden terminar los siguientes enfrentamientos.

\section{Estructura de la memoria}\label{estructura-de-la-memoria}

La estructura que sigue la memoria es la siguiente:

\begin{itemize}
\tightlist
\item
  \textbf{Introducción:} descripción breve del problema que se resuelve y contexto del proyecto junto con la estructura de la memoria y de los anexos.
\item
  \textbf{Objetivos del proyecto:} lista con los objetivos que se quieren conseguir alcanzar con la realización del proyecto. De desglosan en los objetivos generales, técnicos y personales.
\item
  \textbf{Conceptos teóricos:} explicación sencilla sobre los conceptos teóricos que son necesarios para comprender correctamente el proyecto realizado.
\item
  \textbf{Técnicas y herramientas:} listado que explica las técnicas y las herramientas que se han utilizado durante el desarrollo del proyecto.
\item
  \textbf{Aspectos relevantes del desarrollo:} se muestran los aspectos más importantes que han tenido lugar durante el tiempo de desarrollo del proyecto.
\item
  \textbf{Trabajos relacionados:} investigación de trabajos similares o con relación al proyecto desarrollado y comparación de funcionalidades con estos.
\item
  \textbf{Conclusiones y líneas de trabajo futuras:} conclusiones obtenidas tras realizar el proyecto y posibles mejoras que se puedan realizar en actualizaciones futuras.
\end{itemize}



\section{Estructura de los anexos}\label{estructura-anexos}

\begin{itemize}
\tightlist
\item
  \textbf{Plan de proyecto:} descripción de la planificación del proyecto, tiene tanto la planificación temporal como el estudio de la viabilidad económica y legal.
\item
  \textbf{Especificación de requisitos:} se especifican los requisitos en función de los objetivos que se establecieron al inicio del proyecto, además de los que han surgido durante el proyecto.
\item
  \textbf{Diseño:} explicación del diseño que se han seguido para realizar la aplicación que forma el proyecto.
\item
  \textbf{Manual del programador:} explicación de conceptos más específicos y relacionados con la parte de desarrollo y programación del proyecto.
\item
  \textbf{Manual de usuario:} manual que sirve para comprender como se utiliza la aplicación con todas y cada una de las funcionalidades que aporta.
\item
  \textbf{Sostenibilización curricular:} documentación sobre los temas relacionados con la sostenibilidad que se llevan a cabo en el proyecto.
\end{itemize}


\section{Materiales adjuntos}\label{materiales-adjuntos}

Los materiales que se entregan adjuntos con la memoria son: 

\begin{itemize}
\tightlist
\item
	3 Memorias USB con todo lo necesario (memoria, anexos, códigos y documentos explicativos).
\item	
	Vídeo de demostración del funcionamiento. \href{https://youtu.be/CiM1DY7O88A}{Vídeo de demostración de funcionamiento}
\end{itemize}

Además, los siguientes recursos están accesibles a través de internet:

\begin{itemize}
\tightlist
\item
Aplicación web para la liga de tenis de los colegios de Burgos.
    Se puede acceder desde este enlace: \href{http://ligatenisburgos.es/}{Liga de tenis de colegios de Burgos}

\item Repositorio del front-end: 
    \href{https://github.com/Nachoozg/tfg}{LigaTenis front en GitHub}.
    
\item Repositorio del back-end: 
    \href{https://github.com/Nachoozg/tfg-back}{LigaTenis back en GitHub}.

\end{itemize}

\capitulo{2}{Objetivos del proyecto}

En este apartado se van a detallar los objetivos que se quieren alcanzar realizando este proyecto, haciendo una diferenciación entre los objetivos generales, que están relacionados con los requisitos que va a cumplir el software, objetivos técnicos, que son los aprendizajes que necesitaré para alcanzar los objetivos generales y los objetivos personales, relacionados con el punto final al que quiero llegar con este proyecto.

\section{Objetivos generales}\label{objetivos-generales}

\begin{itemize}
\tightlist
\item	
\textbf{Desarrollar la aplicación web completa:} Implementar toda la lógica de programación necesaria para llegar al producto software final.

\item
\textbf{Desarrollar la parte front-end de la web:} La aplicación es fullstack, por tanto, el front-end vendrá desarrollado en Angular que permitirá aprovechar su estructura para que la aplicación sea escalable fácilmente.

\item
\textbf{Desarrollar la parte back-end de la web:} El back-end vendrá desarrollado en .Net, permitiendo desarrollar APIs propias con las que alterar información de bases de datos. Además, el desarrollo de las tablas de la base de datos se llevará a cabo usando HeidiSQL.

\item
\textbf{Sistema de autenticación:} Implementar un sistema que permita validar credenciales de usuarios con diferentes roles que permitirán ver o alterar una información diferente en función del rol asignado. De esta manera, el mantenimiento de los datos será seguro, ya que nadie que no tenga el rol indicado podrá acceder a ella y modificarla.

\item
\textbf{Mostrar información necesaria sobre la liga:} Proporcionar información relevante sobre los colegios que participan en la liga, los jugadores que hay en cada colegio compitiendo y estadísticas de clasificación general, así como información necesaria sobre los partidos jugados y por jugar entre los jugadores representantes de los diferentes colegios.

\item
\textbf{Generar predicciones sobre los partidos por jugar:} Agregar una predicción de resultados de los partidos que se van a jugar más adelante mediante la evaluación de estadísticas de anteriores encuentros entre los colegios y últimos encuentros disputados.
\end{itemize}




\section{Objetivos técnicos}\label{objetivos-tecnicos}

\begin{itemize}
\tightlist
\item
\textbf{Implementar Bootstrap:} para mejorar la apariencia y la forma en la que se navega en la página web. Al implementar esta herramienta, permite que la web sea responsive, es decir, tenga un diseño que se adapte al dispositivo en el que se está viendo y la funcionalidad cambie en los menús y barras de navegación, por ejemplo.

\item
\textbf{Implementar Angular Material:} Permite añadir barras de navegación y menús con una estructura ya definida y que evitan tener que partir de cero al hacer estos componentes.

\item
\textbf{Implementar Toastr:} Permite mostrar alertas que tienen una estructura y permiten ser llamadas en el código de una forma definida, con la que no hace falta crear las tarjetas de alerta desde 0 y hacer su funcionalidad y su diseño.

\item
\textbf{Implementar Full Calendar:} Permite utilizar el diseño y la funcionalidad de un calendario, al que se le pueden añadir eventos y botones como el cambio de mes sin tener que hacer un calendario con sus diferentes funcionalidades desde cero.

\item
\textbf{Implementar Leaflet y OpenStreetMap:} Implementar esta librería (Leaflet) y esta API gratuita (OpenStreetMap) permite mostrar un mapa sobre el cual se puede indicar la ubicación en la que se juega un partido a través de una interacción directa con el mapa que se muestra en la interfaz del usuario.
\end{itemize}


\section{Objetivos personales}\label{objetivos-personales}

\begin{itemize}
\tightlist
\item
\textbf{Mejorar el nivel en el desarrollo del front-end con Angular:} Conseguir mucha experiencia desarrollando aplicaciones web y utilizar al mejor nivel que pueda todas las posibilidades y opciones que ofrece angular respecto a otros frameworks.

\item
\textbf{Mejorar el nivel en el desarrollo del back-end con .Net para el desarrollo de APIs:} Conseguir experiencia y desarrollar conocimientos avanzados sobre el desarrollo de APIs para poder manejar mis propios datos de forma segura y controlada.

\item
\textbf{Ampliar mi capacidad de diseñar interfaces de usuario:} Conseguir mejoras significativas en mi forma de desenvolverme a la hora de crear interfaces de usuario intuitivas y sencillas de usar para el usuario final que vaya a utilizar la aplicación, consiguiendo crear diseños atractivos y útiles.

\item
\textbf{Ampliar mis conocimientos sobre bases de datos:} Mejorar mis conocimientos sobre el uso de bases de datos en este caso mediante el uso de MySQL a través de HeidiSQL que permite crear tablas para bases de datos de una forma más visual, intuitiva y fácil.

\item
\textbf{Aprender a integrar una buena seguridad en los datos que hay detrás de una aplicación:} Conseguir adquirir buenos conocimientos sobre la seguridad de los datos en aplicaciones reales y sobre cómo se lleva a cabo el manejo de ellos sin filtrar datos privados o que ciertos tipos de usuarios no deberían de poder visualizar.
\end{itemize}
\capitulo{3}{Conceptos teóricos}

En esta sección se presentan algunos conceptos teóricos relacionados con el desarrollo de la aplicación.

\section{CRUD}
En el ámbito de la programación un CRUD es un acrónimo que agrupa las cuatro operaciones básicas que son necesarias para gestionar y manejar datos en una aplicación:

\imagen{crud-logo.jpg}{Definición de CRUD. \cite{img:crud-logo}}{0.6}

\subsection{Create (Crear)}
Consiste en la inserción de nuevos registros o nuevos datos en una base de datos. En nuestro caso, dar de alta un usuario o meter un partido nuevo.

\subsection{Read (Leer)}
Consiste en la consulta de los datos ya existentes en la base de datos. Permite mostrar listados con datos, detalles o resultados de búsquedas sin modificar la información original. En nuestra aplicación, consultar la lista de partidos guardados.

\subsection{Update (Actualizar)}
Consiste en hacer modificaciones de registros o datos ya existentes. Se utiliza para corregir, completar o cambiar valores. En este proyecto, actualizar la fecha de un partido guardado.

\subsection{Delete (Eliminar)}
Consiste en borrar registros o datos que ya no son necesarios. Implica borrar los datos de la base de datos. En nuestro caso, borrar un partido que finalmente no se jugará.


Cada operación de un CRUD se puede traducir como sentencias a nivel de base de datos.
Por ejemplo INSERT para Create, SELECT para Read, UPDATE para Update, DELETE para Delete.
 


\section{API}
El concepto de API (Application Programming Interfaces) o Interfaz de Programación de Aplicaciones engloba las funciones, herramientas, reglas y endpoints que nos permiten tener una comunicación directa para poder extraer información o datos, además de funcionalidades interesantes.

Las APIs son muy necesarias para proyectos como este que tienen un CRUD. Las APIs exponen puntos de entrada (endpoints) que permiten ejecutar operaciones CRUD de forma controlada desde puntos externos a la API (en este caso desde la interfaz de la página web). El funcionamiento consiste en asignar a cada endpoint una operacion CRUD de forma clara. Por ejemplo POST /usuarios para Create, GET /usuarios/{id} para Read, PUT /usuarios/{id} para Update, DELETE /usuarios/{id} para Delete.

Las APIs pueden ser abiertas, que son APIs de código abierto que permiten obtener datos gratuitos. Pueden ser internas o propias, creadas por un usuario (como en este caso, creadas por mí), pueden ser de pago, que necesitan que un usuario pague una suscripción para poder sacar datos y usar los endpoints o funcionalidades.


\section{Endpoints}
Los endpoints son los puntos de acceso que se definen dentro de una API y que permiten la comunicación entre aplicaciones. Cada uno de los endpoints corresponden a una URL en específico y tienen asociada una determinada operación como puede ser la obtención de datos, creación de algún recurso o herramienta o actualización de la información. Estos endpoints actúan como puertas de acceso, recibiendo solicitudes del cliente y devolviendo respuestas. Lo más normal es combinarlos con métodos HTTP.


\section{Métodos HTTP}
Los métodos http son las acciones que van a indicar al servidor que tipo de operacion se quiere realizar sobre un determinado recurso cuando se crea una solicitud a través de una API. Los métodos más comunes y los que se manejan en la aplicación creada son:
\subsection{GET}
Solicita datos sin modificarlos.

\subsection{POST}
Crea un nuevo dato o un nuevo recurso.

\subsection{PUT}
Actualiza un dato o recurso que ya existe.

\subsection{DELETE}
Elimina un dato o recurso que ya existe.

Existen otros métodos como el OPTIONS o el PATCH que sirven para consultar capacidades de un endpoint
y para hacer actualizaciones parciales respectivamente.


\section{Front-End}
El front-end es una de las capas de la aplicación web, en concreto la capa con la que interactúa de forma directa el usuario. En este caso está construido con tecnologías como HTML, CSS y TypeScript porque he usado Angular. Se encarga elegir que datos se presentan y como se presentan (colores, estilos, tipos de letra, etc.).

\section{Back-End}
El back-end es la otra capa de la aplicación web, que es la parte que es invisible para el usuario final y se ocupa de llevar toda la lógica, manejo de datos y la seguridad de los datos. Está implementado en este caso en .NET y conecta con base de datos (en este caso MySQL en HeidiSQL) en las que almacena, recupera y actualiza información. Esta parte de la web expone APIs y endpoints que despues consumirá el front-end para mostrar la información.

\section{Full-Stack}
El full-stack es la combinación de las dos capas de la aplicación (front-end y back-end), coordinando la lógica con la representación de los datos. Además, engloba la gestión de las bases de datos y la integración de algunos servicios externos.

\imagen{front-back.png}{Visual Front vs Back. \cite{img:front-back}}{0.7}
\capitulo{4}{Técnicas y herramientas}

\section{Angular}\label{angular}
Angular es un framework de desarrollo de aplicaciones web que es de código abierto y está mantenido por Google, está basado en TypeScript, que es similar a JavaScript y aporta un tipado estático y unas características avanzadas de ES6+, que mejoran la escalabilidad y la capacidad de mantenimiento de los proyectos. 

\imagen{angular-logo.png}{Logotipo de Angular \cite{img:angular-logo}}{0.3}

El principal núcleo que diferencia y da ventaja a Angular respecto a otros frameworks es el concepto de componentes, que son unidades que combinan html, typescript y css. Cada componente representa una de las partes que se ven en la interfaz de usuario y se pueden anidar unos con otros formando jerarquías que definirán la estructura final que tendrá la aplicación. Estos componentes se agrupan en modulos, que permiten organizar el código en diferentes funcionalidades y admitir de esa forma la carga perezosa (lazy loading) para que se cargue solo la informacion esencial en cada ventana y se mejore el rendimiento.

Para hacer la gestión de la comunicación entre componentes y servicios, Angular tiene un sistema que se llama inyección de depencencias, con el que se registran en los injectors o contenedores los servicios que sirven para acceder a datos.

Angular también incluye un enrutador (Router) que permite definir rutas y navegación dentro de la aplicación web, gestionando los parámetros, posibles rutas hijas y los guards (servicios que permiten controlar el acceso a las rutas de la aplicación) que controlan el acceso a apartados de la aplicación según los permisos que se quieran dar. Junto con angular cli (siguiente apartado), ofrece un ecosistema muy maduro. lo que lo convierte en una opción sólida para construir aplicaciones web robustas, mantenibles y de gran escala.

\section{Angular CLI}\label{angular-cli}
Es una herramienta de línea de comandos que acelera la creación, compilación y el despliegue de los proyectos, de manera que permite aplicar mejores prácticas de forma automática.

\imagen{angular-cli.png}{Angular CLI}{0.6}

\section{TypeScript}\label{typescript}
TypeScript es un superset tipado de JavaScript que está desarrollado por Microsoft, lo que indica que todo el código JavaScript válido también lo es en TypeScript, pero con algunas capacidades añadidas como la declaración de tipos estáticos. Está diseñado para desarrollar aplicaciones grandes y capaz de detectar errores a la hora de compilar y evitar problemas.

Con la declaración de tipos estáticos se define una diferencia importante con JavaScript, ya que permite tipificar y que el lenguaje usado sea fuertemente tipado.

\subsection{Ventajas de TypeScript}

\begin{itemize}
  \item \textbf{Buen mantenimiento y escalabilidad}:
        Para proyectos grandes, los tipos ofrecen una red de seguridad que permite evitar errores y facilita el entendimiento de interfaces y unión entre módulos.
  \item \textbf{Mejor experiencia de desarrollo}: 
        Los editores de código (IDEs) pueden ofrecer autocompletado, navegación entre definiciones y verificación de errores gracias al análisis estático.
  \item \textbf{Características avanzadas}:
        Tiene genericos, decoradores y el sistema de tipos que evoluciona para cubrir patrones complejos.
\end{itemize}

\imagen{typescript-logo.png}{Logotipo de TypeScript \cite{img:typescript-logo}}{0.25}

\section{HTML y CSS}\label{html-css}
Estos dos lenguajes de programación se utilizan en conjunto para crear la gran mayoría de páginas web.

\subsubsection{HTML}

HTML (HyperText Markup Language) es un lenguaje de marcado basado en etiquetas que define la parte del contenido que se va a mostrar en una web. Cada elemento como encabezados (\texttt{<h1>–<h6>}), párrafos (\texttt{<p>}), enlaces (\texttt{<a>}) o listas (\texttt{<ul>}, \texttt{<ol>}) se escribe con una etiqueta para abrir y cerrar.

\imagen{html-logo.png}{Logotipo de HTML \cite{img:html-logo}}{0.2}

\subsubsection{CSS}

CSS (Cascading Style Sheets) es un lenguaje de hojas de estilo que indica cómo se deben mostrar los elementos del HTML. A través de selectores y propiedades (por ejemplo, \texttt{color}, \texttt{font-size}, \texttt{margin} o \texttt{display}), CSS permite mantener separados el contenido y su estilo, facilita el mantenimiento y mejora del diseño y permite adaptar la interfaz de usuario a diferentes tamaños de pantalla con técnicas responsive.

\imagen{css-logo.png}{Logotipo de CSS \cite{img:css-logo}}{0.15}

\section{.NET para APIs}\label{net}

.NET, a través de ASP .NET Core, ofrece un framework que es ligero y modular y además está especialmente optimizado para construir APIs RESTful con un alto rendimiento y buena escalabilidad. Su modelo de programación está basado en controladores (\texttt{Controllers}), rutas y servicios.

\begin{itemize} 
  \item \textbf{Documentación automática con Swagger}: integración de forma nativa con \texttt{Swashbuckle.AspNetCore} para generar especificaciones OpenAPI y poder exponer Swagger UI sin tener que crear código adicional.  
  \item \textbf{Entity Framework Core}: ORM (Object-Relational Mapping o Mapeo Objeto-Relacional) ligero y extensible para mapear modelos a bases de datos relacionales, con soporte para migraciones, consultas y seguimiento de cambios.  
  \item \textbf{Seguridad y validación}: tiene soporte para autenticar las APIs basada en tokens como pueden ser JWT y OAuth2.  
\end{itemize}

\imagen{net-logo.png}{Logotipo de .NET Core \cite{img:net-logo}}{0.25}

\section{Visual Studio Code}\label{visual-studio-code}

Visual Studio Code (VS Code) es un editor de código gratuito y multiplataforma que está desarrollado por Microsoft. Combina la ligereza y rapidez con funcionalidades avanzadas típicas de un IDE (entorno de desarrollo integrado), ofreciendo un consumo de recursos moderado en distintas plataformas. Visual Studio Code se centra en mejorar la productividad del desarrollador mediante una interfaz intuitiva y un sistema de pestañas que permite navegar entre ficheros facilmente.

Entre sus características más destacadas están:
\begin{itemize}
  \item \textbf{IntelliSense}: sugerencias de código inteligentes y autocompletado, permitiendo una experiencia mejor y una agilidad mayor. 
  \item \textbf{Depuración integrada}: permite poner puntos de interrupción, inspeccionar variables saltar paso a paso y tiene una consola de depuración unificada.
  \item \textbf{Terminal integrado}: proporciona acceso directo a la línea de comandos dentro del propio editor, permitiendo ejecutar comandos sin tener que salir del editor.  
  \item \textbf{Control de versiones Git}: tiene un panel que enseña las diferencias, historial de commits, ramas y etiquetas y tiene botones para realizar acciones rápidas para commit, push, pull y resolución de conflictos.  
  \item \textbf{Extensiones y Marketplace}: ofrece muchas extensiones para diferentes lenguajes, temas de estilo, depuradores y distintas utilidades.  
  \item \textbf{Workspaces y configuración}: el espacio de trabajo está separado por carpetas y tiene una interfaz que permite la navegación sencilla por ellas.  
\end{itemize}

Este ha sido el editor que he usado para desarrollar la parte del front hecho con Angular, por todas estas ventajas que tiene y sobretodo por las facilidades y comodidades que aportan las extensiones.

\imagen{vscode-logo.png}{Logotipo de Visual Studio Code \cite{img:vscode-logo}}{0.2}

\section{Visual Studio 2022}\label{visual-studio}

Visual Studio 2022 es un entorno de desarrollo integrado (IDE) para Windows, está diseñado para dar soporte a proyectos de gran tamaño y múltiples lenguajes. Este entorno, aprovecha más memoria y potencia de CPU, lo que se traduce en tiempos de carga más rápidos, soluciones con muchos archivos y depuración fluida sin necesidad de subdividir proyectos.

Entre sus características más relevantes están:
\begin{itemize}
  \item \textbf{Plantillas de API ASP.NET Core}: permite crear proyectos Web API con un solo click, incluyendo la estructura de los controladores, archivos \texttt{Program.cs} y \texttt{launchSettings.json} preconfigurados para ejecutar en IIS Express.
  \item \textbf{Integración nativa con Swagger}: una de las características interesantes es que añade y configura \texttt{Swashbuckle.AspNetCore} directamente desde el asistente de NuGet, que genera automáticamente especificaciones OpenAPI y expone Swagger UI para probar los endpoints que se van creando.
  \item \textbf{IntelliSense y refactorización para C\#}: implementa una herramienta que ayuda a completar el código de forma inteligente y con sugerencias basadas en patrones de uso de .NET, que facilitan la creación de controladores, modelos y servicios con interfaces fuertemente tipadas.
  \item \textbf{Extensibilidad y extensiones .NET}: permite instalar extensiones como Pomelo Entity o Entity Framework Tools para extender capacidades de los proyectos de API sin salir del IDE.
\end{itemize}

El motivo principal de selección de este editor para hacer el backend ha sido las facilidades que da para crear para crear proyectos de tipo API NET Core y el uso sencillo de Swagger.


\imagen{visualstudio-logo.png}{Logotipo de Visual Studio 2022 \cite{img:vistualstudio-logo}}{0.2}


\section{Swagger}\label{swagger}

Swagger es un conjunto de herramientas que están basadas en la especificación de OpenAPI que permite diseñar, documentar y probar las APIs de forma automática. En el ecosistema .NET, su integración se realiza habitualmente mediante el paquete NuGet \texttt{Swashbuckle.AspNetCore}, lo que ofrece:

\textbf{Interfaz interactiva (Swagger UI)}: expone una web donde se pueden explorar y ejecutar todos los endpoints de la API sin herramientas externas.

\textbf{Validación y pruebas rápidas}: facilita la comprobación de respuestas y códigos de estado en tiempo real, acelerando el proceso de desarrollo y la detección de errores.


\section{LaTeX}

LaTeX es un sistema de preparación de documentos que está basado en el lenguaje TeX, diseñado para producir documentos una buena calidad. Es especialmente valorado en entornos más académicos y científicos debido a su capacidad para gestionar fácilmente fórmulas matemáticas complejas, bibliografías y diferentes idiomas.

\imagen{latex-logo.png}{Logotipo de LaTeX \cite{img:latex-logo}}{0.3}

\subsection{Overleaf}

Overleaf es un editor de LaTeX basado en la nube que ofrece un entorno colaborativo y accesible desde cualquier navegador. En el, integra un compilador en tiempo real, que permite hacer una vista previa del documento y tiene detección de errores, esto facilita el trabajo en equipo y la revisión en equipos de documentos académicos, artículos y tesis. Permite crear e iniciar proyectos sin tener la necesidad de instalar ningún software localmente y compartir fácilmente el resultado.

\imagen{overleaf-logo.jpg}{Logotipo de Overleaf \cite{img:overleaf-logo}}{0.2}

\section{Control de versiones}\label{control-versiones}
Git es un sistema de control de versiones distribuido que está diseñado para poder gestionar el historial de cambios en proyectos de software.
\href{https://git-scm.com/}{Git}

\imagen{git-logo.png}{Logotipo de Git \cite{img:git-logo}}{0.3}

\section{Github}\label{github}
\href{https://github.com/}{Github} es una plataforma en la nube que permite subir proyectos de software que utilizan el sistema de control de versiones Git. 

\imagen{github-logo.png}{Logotipo de GitHub \cite{img:github-logo}}{0.2}

Subir código a un repositorio de GitHub implica los siguientes pasos básicos:

\subsection{Inicializar o clonar un repositorio Git}
Dentro de la carpeta de un proyecto, permite crear un repositorio local.
Cuando hay un proyecto existente, permite clonarlo para traer una copia local.

\subsection{Gestionar cambios con commits}
Después de crear o modificar archivos, se añaden al área de preparación (staging area) y se registran los cambios con un commit.

\subsection{Conectar el repositorio local con GitHub}
Si no se ha clonado, se puede añadir una URL remota en la que subir el código.

\subsection{Enviar (push) cambios al repositorio en GitHub}
Para subir la rama principal

\subsection{Trabajo colaborativo y ramas}
Ramas (`branches`): consiste en crear nuevas líneas de desarrollo y cambiarse a ellas.
Pull requests (PR): en la interfaz web de GitHub, proponer que tu rama se fusione (“merge”) en la rama principal tras revisión de código.

\subsection{Beneficios de usar GitHub}
Historial de versiones: permite regresar a estados anteriores del proyecto.
Colaboración: múltiples desarrolladores pueden trabajar simultáneamente sin pisarse cambios.
Visibilidad: proyectos open-source se comparten con la comunidad; repositorios privados para código interno.

Con lo que, subir los códigos a GitHub no consiste solamente en copiar y pegar ficheros, sino en aprovechar todo un ecosistema de control de versiones, asegurando que cada cambio quede documentado, pueda revisarse y desplegarse siguiendo buenas prácticas de desarrollo.


\section{HeidiSQL con MySQL}\label{heidi-sql}
HeidiSQL es un cliente gráfico gratuito para Windows diseñado principalmente para gestionar bases de datos MySQL, MariaDB y otros motores compatibles (como por ejemplo PostgreSQL o Microsoft SQL Server). 

\imagen{heidi-logo.png}{Logotipo de HeidiSQL \cite{img:heidi-logo}}{0.25}

Al usarlo con MySQL, HeidiSQL ofrece una interfaz intuitiva que permite:
Conectarse y navegar por servidores MySQL mediante un administrador de sesiones, donde puedes guardar múltiples configuraciones de conexión (host, puerto, usuario, contraseña, etc.).

Visualizar y editar estructuras de bases de datos, tablas, vistas y procedimientos almacenados con un explorador de objetos que muestra de forma jerárquica todos los elementos.
Ejecutar consultas SQL en un editor con resaltado de sintaxis, autocompletado y pestañas múltiples para trabajar en varios scripts a la vez.

Importar y exportar datos de forma sencilla: puedes volcar estructuras y datos a archivos SQL, CSV, XML o incluso sincronizar bases de datos entre servidores con su herramienta de comparación.

Monitorear la actividad del servidor MySQL (procesos activos, status variables) y obtener estadísticas básicas de rendimiento.

\imagen{mysql-logo.png}{Logotipo de MySQL \cite{img:mysql-logo}}{0.4}

HeidiSQL es muy útil cuando se necesita un acceso rápido y eficiente a bases de datos MySQL sin complicaciones de configuraciones avanzadas ni licencias de pago.

\section{Google AI Studio e integración en .NET}\label{google-ai-studio}
Google AI Studio (antes conocido como Vertex AI Workbench) es la plataforma de Google Cloud que unifica herramientas para el desarrollo, entrenamiento, despliegue y monitorización de modelos de inteligencia artificial y machine learning. 

\imagen{ai-studio-logo.png}{Logotipo de Google AI Studio \cite{img:ai-studio-logo}}{0.5}

\subsection{Integración con Google AI Studio en .NET}

La aplicación .NET se conecta a Google AI Studio empleando las librerías oficiales de Google Cloud para .NET, lo que permite aprovechar capacidades avanzadas de IA directamente desde nuestro entorno de desarrollo. El proceso de integración consta de cuatro pasos clave:

\begin{enumerate}
  \item \textbf{Instalación}: añadir el paquete NuGet  
    \texttt{Google.Cloud.AIPlatform.V1} al proyecto.
  \item \textbf{Autenticación}: configurar y utilizar una cuenta de servicio  
    cuyas credenciales se cargan en las variables de entorno.
  \item \textbf{Invocación de la API}: llamar a los métodos de la biblioteca  
    para entrenar, evaluar o desplegar modelos de Machine Learning.
  \item \textbf{Gestión y monitorización}: para procesar las respuestas de la API y supervisar métricas y registros desde el propio código .NET.
\end{enumerate}

Gracias a esta integración, se pueden automatizar tareas de aprendizaje automático sin abandonar todo el ecosistema de .NET, manteniendo una experiencia de desarrollo coherente.


\section{Herramientas integradas en Angular}

\subsection{Angular Animations}
Angular Animations es un módulo oficial que extiende la capacidad de Angular para crear transiciones y efectos dinámicos en la web. Mediante la API basada en funciones como \texttt{trigger}, \texttt{state} y \texttt{transition}, permite definir animaciones de entrada, de salida y cambios de estado en los componentes, además, garantiza un rendimiento bueno en los navegadores más usados actualmente.

\subsection{Angular Material}
Angular Material es la implementación de Material Design para Angular, tiene componentes que vienen preparados y listos para usar directamente como botones, tarjetas, diálogos, menús, tablas, etc. y utilidades del Angular CDK.

\subsection{Bootstrap}
Bootstrap es un framework CSS muy utilizado que, con bibliotecas como \texttt{ng-bootstrap} o \texttt{ngx-bootstrap}, se integra perfectamente en proyectos Angular. Esto proporciona un sistema con utilidades, estilos y componentes interactivos como modales, barras de navegación y botones, permitiendo acelerar el desarrollo de interfaces de usuario.

\subsection{Leaflet}
Leaflet es una librería JavaScript que sirve para tener en la web mapas interactivos, su integración en Angular se realiza con paquetes como \texttt{@asymmetrik/ngx-leaflet}. Tiene componentes y directivas que facilitan la creación de mapas, edición de mapas y permite poner punteros en los mapas (lo que se ha usado en este caso) para indicar lugares.

\subsection{Toastr}
Toastr es una librería que permite sacar notificaciones de tipo “toast” que, mediante el paquete \texttt{ngx-toastr}, se adapta para mostrar mensajes breves que salen para advertir de diferentes acciones. Permite configurar la posición, el tiempo que se muestran y los estilos de cada notificación a través de servicios, mejorando así la experiencia de usuario.

\subsection{FullCalendar}
FullCalendar es un plugin para añadir calendarios interactivos en la web. La integración con Angular se realiza a través de \texttt{@fullcalendar/angular}, que lo que hace es exponer componentes para renderizar vistas de meses, semanas o incluso días, permite gestionar eventos y personalizar la apariencia.


\section{Draw.io}\label{drawio}

Draw.io es una web gratuita que permite crear y editar diferentes gráficos como diagramas o flujogramas. 

Esta herramienta se usa para generar los diagramas que se muestran en la parte de los anexos.


\section{AWS (Amazon Web Services)}\label{aws}
AWS (Amazon Web Services) es la plataforma de servicios en la nube de Amazon. Ofrece diferentes herramientas para que los desarrolladores desplieguen aplicaciones sin tener que tener ejecutando en local todo. Dentro de este servicio, para este proyecto se utiliza la máquina EC2, que permite crear y ejecutar instancias en la nube. Se usan para alojar sitios web y aplicaciones en producción en un servidor que es escalable.

\imagen{aws-logo.png}{Logotipo de Amazon Web Services \cite{img:aws-logo}}{0.2}

\section{SonarQube}\label{sonarqube}
SonarQube es una plataforma de análisis de código que inspecciona los proyectos y busca errores, vulnerabilidades y errores de mantenibilidad. Ayuda a ver la calidad del software que se está desarrollando y permite detectar problemas antes de finalizar el desarrollo.

\imagen{sonarqube-logo.png}{Logotipo de SonarQube \cite{img:sonarqube-logo}}{0.4}

\section{Gemini}\label{gemini}
Gemini es una familia de modelos de inteligencia artificial (IA) de Google que fue diseñada para ser multimodal, permitiendo procesar tanto texto, como imágenes, audio y código. Fue desarrollado para ser conversacional y generativo, lo que le permite entender y crear contenido nuevo.

En este caso se ha utilizado para generar las imágenes de los jugadores y de los colegios inscritos en la liga.

\imagen{gemini-logo.jpg}{Logotipo de Gemini \cite{img:gemini-logo}}{0.4}
\capitulo{5}{Aspectos relevantes del desarrollo del proyecto}

En este apartado se van a presentar los aspectos más importantes de todo el desarrollo del proyecto. Se incluye desde el inicio del proyecto y la formación realizada hasta todo el proceso de desarrollo del software, el diseño y la resolución de problemas que se han ido planteando con el tiempo.

\section{Inicio del proyecto}\label{inicio-del-proyecto}

La idea de este proyecto apareció como la idea de combinar uno de los deportes que más he practicado y más disfruto viéndolo con mis estudios y la programación.

En concreto, escogí realizar una web para la gestión de una liga de tenis para los colegios de Burgos porque es algo que no existe y es novedoso. En las instalaciones de la escuela en la que iba yo a jugar a tenis había tablones en las paredes en los que se colgaban papeles con tablas de jugadores que estaban inscritos en torneos que se jugarían más adelante. Todo estaba gestionado con los tablones de anuncios y si quería consultar cuando tenía que jugar un partido o quienes estaban inscritos para jugar el torneo, tenía que acercarme hasta el lugar a consultar el papel.

Además, en el caso de el colegio al que yo iba, no había clases de tenis y por tanto, mucho menos había competición o partidos contra otras personas. 

Esa necesidad que vi, junto con el interés que me despertó el desarrollo web cuando lo descubrí, me hicieron que utilizar las herramientas que he usado para el desarrollo (Angular y .NET) me resultara más interesante y más llevadero para aprenderlo.

Una vez llegado a la idea y tras recibir el visto bueno del tutor, me puse con el desarrollo de todo el proyecto.

\section{Formación}\label{formacion}

Las herramientas que se utilizan para desarrollar este proyecto no son herramientas que se hayan visto en las asignaturas de la carrera o se ha tocado muy poco o sin relación al desarrollo web, por tanto, se requerían conocimientos de los que no disponía al comienzo del desarrollo. Para adquirir esos conocimientos que no tenía me puse a investigar en profundidad acerca de estas tecnologías y tuve que realizar cursos formativos con mucho material que me fuera útil.

Además, con el auge de la inteligencia artificial en los últimos años, finalmente surgió la idea de añadir un asistente inteligente que aparecería en forma de chat en la web y también implicaba investigación y adquirir nuevos conocimientos sobre como llevar la idea a cabo.

\subsection{Angular}

Angular era nuevo para mí, sabía que existían frameworks como React y Django pero no Angular. Realicé un par de cursos que me fueron útiles para poder aprender mucho sobre el tema:

\begin{itemize}
\tightlist
\item
  \emph{Legacy - Angular: De cero a experto} (Udemy)
  \cite{course:angular_legacy}.
\item
  \emph{Aprende ANGULAR 17 desde cero para principiantes GRATIS} (Youtube)
  \cite{course:angular_midudev}.
\item
  \emph{Documentación de Angular en español}
  \cite{web:angular_hispano}.
\item
  \emph{Documentación oficial de Angular en inglés}
  \cite{web:angular_oficial}.
\end{itemize}


\subsection{.NET}

Sobre .NET ya tenía algo de conocimiento y además, el lenguaje es similar a Java, que si que se estudia en la carrera y por tanto, partía de algo de experiencia utilizándolo. Aún así, también consulté algunos tutoriales para conseguir obtener más información y experiencia con esta herramienta.

\begin{itemize}
\tightlist
\item
  \emph{Curso .NET Core} (Youtube)
  \cite{web:curso-net}.
\item
  \emph{.NET 8 Web API and Entity Framework} (Youtube)
  \cite{web:net-web-api}.
\item
  \emph{Paseo por el lenguaje C sharp}
  \cite{web:net_tutorial}.
\item
  \emph{Documentación de .NET}
  \cite{web:net_documentacion}.
\item
  \emph{documentación de ASP.NET}
  \cite{web:net-asp}.
\end{itemize}


Además, para todos los problemas que han ido surgiendo a lo largo del desarrollo, también ha sido muy útil recurrir a páginas como \href{https://stackoverflow.com/questions}{StackOverflow} y \href{https://github.com/community}{Github Community}, donde se muestran los mismos problemas que han tenido otros usuarios y otros usuarios responden con las posibles respuestas a los problemas planteados.

\section{Desarrollo}\label{desarrollo}

En esta sección se van a describir los aspectos más relevantes de todo el proceso de desarrollo del proyecto. Todos los aspectos señalados han sido puntos clave para poder conseguir este producto final.

\subsection{Creación del front-end de la web}

Lo primero que se hizo al comenzar con el desarrollo de este trabajo fue crear el proyecto en Angular y comenzar el desarrollo de toda la estructura básica para poder tener el esqueleto de la web creado, con esto ya se podían ver las secciones principales del proyecto en un inicio, la barra de navegación entre las diferentes partes de la web y toda la estructura de formularios y botones conectados entre ellos para después poder comenzar con el desarrollo de la API en back-end y poder hacer operaciones sobre la base de datos.

\imagen{listadoInicial.png}{Estructura inicial sin conexión a base de datos}{0.8}

\imagen{formInicial.png}{Formulario creado inicialmente (Vista del formulario de colegios)}{1.0}

\subsection{Creación de la base de datos}

Antes de pasar a crear la API se creó la base de datos en HeidiSQL con sus diferentes tablas que se necesitarían en un principio y que después con el tiempo fueron evolucionando a tener más columnas o tablas nuevas para almacenar datos nuevos. Con la creación anterior de los formularios al hacer la estructura básica de la web en Angular, ya se tenía el conocimiento necesario para saber que datos eran necesarios para poder tener todo conectado y poder almacenar toda la información que sería necesaria para llevar la correcta gestión de la liga.


\section{Integración de APIs}\label{integracion-apis}
\subsection{Creación e integración de APIs propias}

Uno de los puntos principales del desarrollo fue la creación de la API propia para poder sacar, meter, modificar y consultar datos necesarios de las tablas de la base de datos y así poder tener las tablas para ver la información y los formularios para editar y añadir esa información. Esta API se creó en .NET con el tipo de solución que se ve en la imagen.

\imagen{tipo-proyecto.jpg}{Tipo de proyecto elegido}{0.7}

A medida que se iban creando diferentes endpoints, se iban probando en Swagger para ver si funcionaban correctamente o si de lo contrario fallaba algo y había que arreglarlo. Cuando esos endpoints estaban creados y probados, se creaba la conexión del back con el front a través de los endpoints que se habían creado.

\subsection{Integración de APIs externas}

Otro aspecto relevante fue la integracion del proyecto con APIs externas, que en este caso se utilizaron dos y fueron:

\subsubsection{Leaflet}

Para poder añadir una modal en la web en la que se mostraría un mapa interactivo para ver donde se jugarían los partidos. En el caso de ser un usuario normal, mostraría la ubicación del partido únicamente, pero en caso de ser un usuario con roles de administrador o arbitro, permitiría poner el puntero en el mapa para marcar en que lugar se disputaría el partido que se estaba creando o editando.

\imagen{c1.png}{Ejemplo del mapa que muestra Leaflet}{0.8}

\subsubsection{Google AI}

La integración con esta API permitió añadir un chatbot inteligente a la web para poder hacerle consultas sobre temas relacionados con la liga de tenis y no tener que ir manualmente a consultar las diferentes tablas que estaban en las secciones específicas. Para seleccionar esta API, se eligió entre otras (como por ejemplo, la de OpenAI) porque a la hora de usar la web de forma real, en caso de que se utilizara para gestionar una liga real de tenis para colegios, estas APIs tienen unos costes por consulta y por número de tokens usados en entornos de producción, y google ai studio es la que ofrece menores costes y mejor desempeño y facilidades a la hora de integrarlo en nuestras webs propias.


\section{Diseño}\label{diseño}

El diseño de la interfaz de usuario de la web, también fue uno de los aspectos clave para poder tener una web amigable y que fuera sobre todo intuitiva y fácil de utilizar para el usuario final. Para ello, en un inicio, el diseño era muy simple y se tenía sin colores específicos y con una distribución de los elementos básica solamente con el diseño por defecto al ir creando los componentes y servicios.

\imagen{form-jugador.png}{Formulario inicial para agregar jugadores}{1.0}

Cuando se llegó al punto en el que todo el front estaba perfectamente integrado con el back y se podían hacer todas las operaciones necesarias para gestionar los datos de la web, se comenzó a mejorar la interfaz de usuario añadiendo estilos, colores, tipografía y creando una distribución de los elementos, botones, tablas, imágenes, etc. que permitiera facilitar lo máximo posible el uso de la web.

Estas imágenes muestran algunas ventanas de la aplicación antes y después de empezar con las mejoras de diseño.

\imagen{vista-jugador.png}{Vista de la información del jugador antes de comenzar la mejora de diseño}{0.6}

\imagen{tabla-clasificacion.png}{Vista de la tabla de clasificación sin mejoras en su diseño}{0.6}

\imagen{calendario-inicial.png}{Vista del calendario creado al mejorar el diseño}{0.5}

\imagen{lista-colegios.png}{Listado de colegios al mejorar el diseño}{0.6}

Para ello fue importante ponerse en la posición del usuario final y navegar por la web sin saber como se utilizaba para llegar a conclusiones de las que se sacaría un rediseño nuevo para llegar al producto final. Para lograr este diseño final, se han utilizado varias herramientas y tecnologías de diseño como Bootstrap y Angular Material. Con Bootstrap se ha hecho una gestión de los estilos y creación de botones y barra de navegación. Con Angular Material también se han creado componentes de la web como botones y tablas. Para poder mostrar alertas al usuario y darle advertencias sobre los datos que estaba añadiendo, editando o eliminando, se utilizó toastr, otra librería que permite mostrar alertas emergentes con información específica y con una duración específica.
En los puntos finales del desarrollo, también se ha adaptado el diseño para el uso de la web en dispositivos móviles, haciendo que el diseño sea responsive y se compriman los elementos al tamaño de la pantalla utilizada.

En las imagenes siguientes se muestra el diseño final de las partes de la web mostradas anteriormente en fases de desarrollo:


\imagen{disenoInicio.png}{Diseño final de la página principal de la aplicación}{1.0}

\imagen{vistaJugadorFinal.png}{Vista final de la información del jugador}{1.0}

\imagen{tablaClasFinal.png}{Vista de la tabla de clasificación final}{1.0}

\imagen{listadoColFinal.png}{Listado de colegios final}{1.0}

\imagen{vista-chat.png}{Diseño del chatbot inteligente integrado}{0.3}

Para la creación de las imágenes como el logotipo de la aplicación se ha utilizado la herramienta paint, que permite hacer imágenes básicas como son las de la imagen de la ventana de la web y la pestaña de inicio.
Las imágenes que se han utilizado en la web son:
\begin{itemize}
\tightlist
\item
  \emph{Imágen raqueta de tenis} (Unsplash)
  \cite{img:raqueta-img}.
\item
  \emph{Imágen de hombre jugando a tenis} (Unsplash)
  \cite{img:hombre-tenis}.
\end{itemize}

\section{Resolución de problemas}\label{resolucion-problemas}

Durante el desarrollo del proyecto se han ido presentando varios problemas o desafíos que se han tenido que ir resolviendo para lograr el funcionamiento correcto de todo.

\subsection{Problema de conexión del front-end con el back-end}

El primer problema que surgió fue como hacer bien la conexión entre Angular y .NET para poder manejar los datos de las tablas de la base de datos y hacer las conexiones correctamente. 

Para resolver el problema se tuvieron que crear servicios de Angular que apuntaran a las url generadas por los endpoints que se habían creado en .NET.

\imagen{url-endpoints.png}{Ejemplo de url de conexión con los endpoints desde un servicio en Angular}{0.8}

\subsection{Problema con la subida y visualización de imágenes}

La integración de la posibilidad de meter imágenes en los formularios para los jugadores y los colegios hizo que surgiera el problema de que al subir una imagen, se paraba la ejecución de todo el proyecto automáticamente y daba errores. Ese problema venía porque se estaban metiendo directamente las imágenes a la base de datos y esto hacía que colapsara todo.

Para resolver este problema lo que se hizo fue hacer modificaciones en .NET y configurarlo de forma que las imágenes que se subieran se guardarían en un servidor interno de.NET (wwwroot) y lo que se guardaría en base de datos era solo el enlace a esa imágen para evitar subir mucho peso a la base de datos.


\imagen{wwwroot.png}{Estructura de wwwroot en .NET}{0.6}

\subsection{Problema con los estilos de FullCalendar}

Estilizar el calendario no fue fácil ya que FullCalendar presenta su diseño por defecto y hay configuraciones para cambiar un poco los estilos de los botones y cambiar el idioma por defecto del calendario que es el inglés.

Para dar con la solución se tuvo que investigar durante un tiempo en foros y páginas que aportaban distintas soluciones e ir adaptando poco a poco el diseño para conseguir lo que se quería.

\subsection{Problema con la subida del proyecto a Amazon Web Services}

Antes de hacer este proyecto partía sin idea de como se podía hacer para que la web funcionara directamente buscando su url en cualquier dispositivo en vez de ejecutar la aplicación en local desde mi ordenador.

Para resolver esto y poder cumplir con ello, este vídeo sirvió de ayuda \href{https://www.youtube.com/watch?v=nQdyiK7-VlQ}{Easily deploy full stack app} junto con muchas consultas a foros como StackOverflow y Github Community, donde se pueden ver preguntas similares que hayan hecho otros usuarios antes con los mismos problemas.

Finalmente se escogió una máquina EC2 de AWS (Amazon Web Services) en la que se subieron las tres partes de la aplicación web (proyecto Angular, proyecto .NET y base de datos) para ejecutarlas y dejar la máquina corriendo para mantener la aplicación en ejecución constantemente.

Esta subida también implicó modificar el código, ya que surgieron problemas nuevos porque las url de conexión del front-end con los endpoints generados por el back-end también cambiaron al no ser el localhost el que estaba en ejecución.

\imagen{urlNueva.png}{Cambio de las url de conexión con los endpoints}{0.9}

\section{Análisis del código con SonarQube}

Se ha analizado el código para ver si tiene problemas de seguridad, de fiabilidad, mantenibilidad y ver si hay código duplicado.

Como se puede ver en la imagen, el nivel de seguridad lo ha valorado como A, que es el mejor valor posible, el 0 indica que ha habido 0 vulnerabilidades. En la fiabilidad, tiene el nivel B, que indica que ha encontrado 6 vulnerabilidades con algunos posibles bugs que no afectan en el funcionamiento de la aplicación. En cuanto a la mantenibilidad, la nota también ha sido la mejor posible y indica que puede haber algunos problemas muy leves que no implican una corrección inmediata. El porcentaje de código duplicado también ha salido muy bajo (0.2), lo que indica que el código no se repite de forma innecesaria.

Por otro lado, se ha analizado también el código del back-end, pero en este caso, no hay una interfaz tan visual y es por comandos. La salida ha mostrado algunas advertencias sobre las versiones instaladas de los paquetes que se están utilizando en el proyecto, avisando de que tiene versiones más modernas y con menos problemas de seguridad.

\imagen{sonarqube.png}{Imágen del análisis del código en SonarQube}{0.9}
\capitulo{6}{Trabajos relacionados}

Este apartado sería parecido a un estado del arte de una tesis o tesina. En un trabajo final grado no parece obligada su presencia, aunque se puede dejar a juicio del tutor el incluir un pequeño resumen comentado de los trabajos y proyectos ya realizados en el campo del proyecto en curso. 

\capitulo{7}{Conclusiones y Líneas de trabajo futuras}

Todo proyecto debe incluir las conclusiones que se derivan de su desarrollo. Éstas pueden ser de diferente índole, dependiendo de la tipología del proyecto, pero normalmente van a estar presentes un conjunto de conclusiones relacionadas con los resultados del proyecto y un conjunto de conclusiones técnicas. 
Además, resulta muy útil realizar un informe crítico indicando cómo se puede mejorar el proyecto, o cómo se puede continuar trabajando en la línea del proyecto realizado. 



\printbibliography
\end{document}
